\chapter{Problem specific Routines} \label{sec3}

\medskip

{\bf General remarks}

\medskip

As mentioned earlier, already the hard wired options in EIRENE allow
a large variety of linear and also of non-linear transport problems to be studied. In oder versions of EIRENE (2001 and older) in most
cases of linear neutral gas transport in a prescribed background plasma
the only problem specific FORTRAN that is required consisted of
a few parameter statements (i.e., the COMMON deck PARMUSR described
below in section \ref{sec3.1}).  

In it one had defined the
storage requirements for this particular case.
\\
{\bf In versions EIRENE$_{2002}$ and younger this pre-assignment
of storage has been removed, by a dynamical allocation of storage
throughout.}
\\
Some controllable storage options (e.g.: trading CPU - vs. storage -
optimization), formerly under: ``parameters for storage
reductions'' in PARMUSR, can now be set in an additional, but
optional input line at the beginning of input block 1 (as new
second line in this block), see section \ref{sec2.1}.

\vspace{1em}

The names of all other routines in the user specified block
end with ...USR.

\smallskip

The description of the ``additional tally" routines UPTUSR, UPCUSR,
UPSUSR and UPNUSR given below (section \ref{sec3.2}) is meant as a
guide for the more experienced users of the code only. These
routines allow to extend the number of responses estimated by EIRENE
almost arbitrarily.

The number of options for fast particle surface interaction models
can be increased by adding subroutines REFUSR and SPTUSR,
which have to provide the particle data (after reflection or sputtering, respectively)
(velocity vector, weight, type, species, etc.) after
a test particle has collided with a non-transparent surface and neither
absorption, nor specular reflection nor the ``thermal particle re-emission"
is identified as
surface event by the flags and the random sampling procedure. The use of
this routine REFUSR is described in section \ref{sec3.3}.

By including a birth point sampling routine SAMUSR, one can easily
model any spatial distribution of the primary source, in addition to the
preprogrammed options described in section \ref{sec2.7}.
This is explained in section \ref{sec3.4}.

Some geometrical variables, such as cell volumes, surface coefficients,
etc., may be re-defined or modified after reading the input data.
This is done by a call to routine GEOUSR, called from EIRENE
subroutine INPUT. See section \ref{sec3.5.1} below. Similarly,
any background data (plasma profiles) may explicitly be modified
in routine PLAUSR (see section \ref{sec3.5.2}).

The routine for additional background profiles PROUSR (called for
those background tallies, for which INDPRO=5, see block 5, is described
in section \ref{sec3.6}. It is recommended to use this option,
if a particular closed form expression for the background profile
is needed. It is then more convenient than the external data file
options INDPRO=6, or INDPRO=7, see section \ref{sec4}.

The various routines needed for the general user supplied geometry
option (NLGEN, LEVGEO=10, see block 2a) are described in section
\ref{sec3.8}. Simple examples are also given there.
\newpage

\section{Parameter Statements (for EIRENE-2001 or older)} \label{sec3.1}
The deck ``PARMUSR" is the first part of the user
supplied Fortran
file. It contains parameter statements used in an EIRENE run
and it determines the storage required to run EIRENE on a particular
problem.

In EIRENE versions after 2001 dynamic allocation of storage is implemented, the module
PARMUSR has been removed.
Rather than from fixed ``parameter statements", the storage (such as no. of species, grid size, no. of atomic processes, etc.) is now allocated directly
as identified from input flags, usually with the same name as previously in module PARMUSR.
The actual ``physical" numbers used in a particular run have an additional letter ``I'' at the end of the variabel name.
E.g.  NATM  (storage for number of different atomic species)
then is identical with input flag NATMI (input block 4a),
etc. Only very few exceptions exist,
in which the storage parameters may be larger than the actual physical parameters in a run,
such as: NSTRA and NSTRAI, the number of strata.

Such deviations may happen if a parameter is altered in the code after the automatic storage
parameter evaluation in routine FIND\_PARAM.f was carried out, e.g. in the above example, if the census
stratum is turned off later in the run due to other parameter settings.

The module PARMUSR reads: (... stands for an integer not less than 1)
\begin{verbatim}
*COMDECK PARMUSR
C
C  Geometry
      PARAMETER (N1ST=..., N2ND=..., N3RD=...)
      PARAMETER (NADD=..., NTOR=...)
      PARAMETER (NLIM=..., NSTS=...)
      PARAMETER (NPLG=..., NPPART=...)
      PARAMETER (NKNOT=..., NTRI=...)
      PARAMETER (NCOORD=..., NTETRA=...)
C  Primary Source
      PARAMETER (NSTRA=..., NSRFS=...)
      PARAMETER (NSTEP=...)
C  Species and Tallies
      PARAMETER (NATM=..., NMOL=..., NION=..., NPLS=...,
                 NADV=..., NADS=...)
      PARAMETER (NCLV=..., NSNV=..., NALV=..., NALS=..., NAIN=...)
      PARAMETER (NCOP=..., NBGK=...)
C  Statistics
      PARAMETER (NSD=..., NSDW=..., NCV=...)
C  Atomic Data
      PARAMETER (NREAC=..., NRRC=..., NREI=...)
      PARAMETER (           NRCX=..., NREL=..., NRPI=...)
C  Surface Reflection Data
      PARAMETER (NHD1=..., NHD2=..., NHD3=..., NHD4=...,
                 NHD5=..., NHD6=...)
C  e.g.: TRIM Database
C     PARAMETER (NHD1=12, NHD2=7, NHD3=5, NHD4=5, NHD5=5, NHD6=...)
C  Diagnostic Chords Data
      PARAMETER (NCHOR=..., NCHEN=...)
C  Interfacing routines:
      PARAMETER (NDX=..., NDY=..., NFL=...)
C  Census Arrays
      PARAMETER (NPRNL=...)
C
C  PARAMETERS FOR STORAGE REDUCTIONS
C
C  1.) IGJUM-FLAGS FOR SPEEDUP OF GEOMETRICAL CALCULATIONS
C
C  NOPTIM: REDUCES IGJUM3-ARRAY: IGJUM3(NOPTIM,NLIMPS)
C          NOPTIM=N1ST*N2ND*N3RD+NADD:  NO STORAGE OPTIMIZATION,
C                                       SPEEDUP OF GEOMETRICAL
C                                       CALCULATIONS IS POSSIBLE
C                                       BY CH3 FLAGS
C          NOPTIM=1 MINIMAL STORAGE, NO SPEEDUP OF
C                                    GEOMETRICAL CALCULATION
C  DEFAULT:
      PARAMETER (NOPTIM=...)
C
C  NOPTM1: BIT-ARITHMETIC/INTEGER, E.G. 32 FOR IBM RISK, 46 FOR CRAY
C          DEFAULT: NOPTM1=1: NO STORAGE OPTIMIZATION
C  DEFAULT:
      PARAMETER (NOPTM1=...)
C
C  2.) EXTERNAL GEOMETRY?
C  NGEOM_USR = 1  ==> EXTERNAL GEOMETRY ROUTINES (...USR), LEVGEO=10
C                     NO STORAGE FOR GEOMETRY DATA IN EIRENE
C  NGEOM_USR = 0  ==> ELSE, NO STORAGE OPTIMIZATION
C  DEFAULT:       PARAMETER (NGEOM_USR=...)
C
C  3.) SUM OVER STRATA
C  NSMSTRA = 0  ==> SUM OVER STRATA IS NOT PERFORMED
C  NSMSTRA = 1  ==> SUM OVER STRATA IS PERFORMED
C  DEFAULT:
      PARAMETER (NSMSTRA=...)
C
C  4.) CALCULATION OR STORAGE OF ATOMIC DATA
C
C  NSTORAM=0,1,2...,8,9.    =0: MINIMUM STORAGE, MAXIMUM CALCULATION
C                           =9: MAXIMUM STORAGE, MINIMUM CALCULATION
C  DEFAULT:
      PARAMETER (NSTORAM=...)
C
C  5.) SPATIALLY RESOLVED SURFACE TALLIES?
C
C  NGSTAL = 0  ==> NO STORAGE FOR SPATIALLY RESOLVED SURFACE TALLIES
C  NGSTAL = 1  ==> SPATIALLY RESOLVED SURFACE TALLIES ARE COMPUTED
C  DEFAULT:
      PARAMETER (NGSTAL=...)
\end{verbatim}
{\bf Meaning of the Parameter Variables (and Defaults)}
\begin{list}{ }{}
\item[{\bf N1ST }]  Maximum number of standard mesh points in x-
(radial) direction (block 2a)
\item[{\bf N2ND }] Maximum number of standard mesh points in y-
(poloidal) direction (block 2b)
\item[{\bf N3RD }] Maximum number of standard mesh points in z-
(toroidal) direction (block 2c)
\item[{\bf NADD }] Maximum number of additional zones defined by the
additional surfaces (block 2d)
\item[{\bf NTOR }] Maximum number of toroidal segments for
approximation of torus by cylindrical segments (block 2c)
\item[{\bf NLIM }] Maximum number of ``additional surfaces" (block 3b)
\item[{\bf NSTS }] Maximum number of non-default standard surface
models (block 3a)
\item[{\bf NPLG }] Maximum number of points in one polygon (block 2a)
\item[{\bf NPPART }] Maximum number of valid parts in one polygon (block 2a)
\item[{\bf NKNOT }] Maximum number of knots for mesh of triangles (block 2a)
\item[{\bf NTRI }] Maximum number of triangles (block 2a)
\item[{\bf NCOORD }] Maximum number of knots for mesh of
tetrahedrons (block 2a)
\item[{\bf NTETRA }] Maximum number of tetrahedrons (block 2a)
\item[{\bf NSTRA }] Maximum number of strata (sources) (block 7)
\item[{\bf NSRFS }] Maximum number of sub-strata in one stratum (block 7)
\item[{\bf NSTEP }] Maximum number of different step functions used for
primary source sampling (block 7)
\item[{\bf NATM }] Maximum number of different atom species (block 4a)
\item[{\bf NMOL }] Maximum number of different molecule species (block 4b)
\item[{\bf NION }] Maximum number of different test ion species (block 4c)
\item[{\bf NPLS }] Maximum number of different bulk ion species (block 5)
\item[{\bf NADV }] Maximum number of additional volume averaged
tallies, track-length estimated (block 10a)
\item[{\bf NADS }] Maximum number of additional surface averaged
tallies  (block 10d)
\item[{\bf NCLV }] Maximum number of additional volume averaged
tallies, collision estimated (block 10b)
\item[{\bf NSNV }] Maximum number of additional volume averaged
tallies, snapshot estimated (in time dependent mode only) (block 13)
\item[{\bf NALV }] Maximum number of algebraic
tallies obtained from volume
averaged tallies (block 10c)
\item[{\bf NALS }] Maximum number of algebraic
tallies obtained from
surface averaged tallies (block 10e)
\item[{\bf NAIN }] Maximum number of additional input tallies, not
needed by EIRENE, but to facilitate presentation (block 14)
\item[{\bf NCOP }] Maximum number of tallies for coupling to external
codes (block 14)
\item[{\bf NBGK }] Maximum number of tallies for iterative BGK scheme
for neutral-neutral interactions, see MODUSR, below.
\item[{\bf NSD }] Maximum number of standard deviation profiles for
volume averaged tallies (block 9)
\item[{\bf NSDW }] Maximum number of standard deviation profiles
for surface averaged tallies (block 9)
\item[{\bf NCV }] Maximum number of covariance estimates (block 9)
\item[{\bf NREAC }] Maximum number of atomic reactions from external
files (block 4)
\item[{\bf NRRC }] Maximum number of re-combination processes (blocks 4,5)
\item[{\bf NREI }] Maximum number of electron impact processes
(blocks 4,5)
\item[{\bf NRCX }] Maximum number of charge exchange reactions (blocks 4,5)
\item[{\bf NREL }] Maximum number of elastic reactions
(blocks 4,5)
\item[{\bf NRPI }] Maximum number of ion impact reactions (blocks 4,5)
\end{list}
Reflection database parameters: see example above. NHD6 is the maximum
number of different target-projectile combinations, for which databases can be included.
\begin{list}{ }{}
\item[{\bf NCHOR }] Maximum number of lines of sight for diagnostics module (block 12)
\item[{\bf NCHEN }] Maximum number of energy values for spectra computed
in diagnostics module (block 12)
\item[{\bf NDX, NDY, NFL }]~

Problem specific, in case of coupling to external code.
See input block 14. Otherwise (stand alone mode): just set to 1.
\item[{\bf NPRNL }] Maximum number of particles stored on census array, for
time dependent mode (block 13)
\end{list}

In addition to these parameters, which depend upon the size of the
particular problem under investigation, there are a few further
parameters in PARMUSR, which can be used to optimize storage versus CPU
performance. E.g., atomic data can either be stored in great detail, or
they can be recomputed whenever they are actually needed. This is controlled
by the parameter NSTORAM.
The latter
choice (e.g., NSTORAM=0, or NSTORAM=1)
may be the better one on very large meshes (above, say, 20000
cells, as are commonly encountered in 3D Stellarator applications).

\newpage
\section{The ``Additional Tally" routines UPTUSR, UPCUSR, UPSUSR and UPNUSR}
\label{sec3.2}

Let $   f(\vr , \vv ,j_S)$ be the distribution function in the 6
dimensional $\mu$ -space for particle species $j_S$, S = PH:
Photons, S = A: Atoms, S = M: Molecules, S = I: Test Ions, S = P:
Bulk Ions

e.g.:
\begin{eqnarray}
& & \\ j_A    & =  & H  , D , T , He , C , O  ,  \dots \nonumber \\
j_M  &  = & H_2  , D_2,  HD , CH_4 , H_2 O ,  \dots \nonumber \\
j_I  &  = & H_2^+ , CH_3^+ ,  C^+ , O^{++} , \dots  \nonumber \\
j_P  &  = & H^+  , D^+ , He^+  ,  He^{++} , C^+ , \dots \nonumber
\end{eqnarray}
The distinction between ``Test Ions" and ``Bulk Ions" is problem specific,
see input blocks 4 and 5.
The ``volume averaged tallies" estimated by EIRENE (see table \ref{tab2})
are profiles of
``responses"
\[
R_g   =    <    f | g   >
=   \sum_{j}   \int d^3 \vv
\int_{volume} d^3 \vr \ g(\vr , \vv
, j)     \cdot     f(\vr , \vv , j)\label{est}
\]
for several preprogrammed and an arbitrary number of additional user
supplied ``detector functions" $g$. Here the summation over $j$ is only
over test particle species, i.e., atoms, molecules or test ions, and
not over background particles (``bulk ions", electrons). The default EIRENE
estimators are not based upon moments of distribution function $f$
but upon moments of the flux
\begin{displaymath}
\Phi ( j_S , \vr ,
\vv )   =   |
\vv |  f(j_S , \vr , \vv),
\end{displaymath}
hence
on detector functions $g^{*}  =   g / | \vv |
$.

Occasionally also
moments of the pre-collision density
\begin{displaymath}
\Psi ( j_S ,
\vr , \vv )   =
\Phi ( j_S , \vr , \vv )
    \cdot     \Sigma ( j_S , \vr , \vv )
\end{displaymath}
are utilized, then
with detector functions $ \hat{g}   =   g /
(| \vv | \ \Sigma)    $.
Here $\Sigma$ (dimension: 1/length) is the local total macroscopic
cross section, i.e., the inverse of the local mean free path length.
The estimators may, in general, be composed of
\\ a ``track-length estimate"
\begin{equation} \label{tracklest}
^k{R_t}  =
\sum_{i=0}^{k_{n-1}}   \ ^k w_i^* \  I_{l,i}
  =   \sum_{i=0}^{k_{n-1}} \ ^k w_i^*
\int_{\vr_i}^{\vr_{i+1}}   dl \  g^{*} (l)     ,
\end{equation}
and a ``collision estimate" :
\begin{equation} \label{collest}
^k{R_c}  =   \sum_{i=1}^{k_n}   \ ^k{\hat{w_i}} \
\hat{g} ( \vr_i )     .
\end{equation}
Here $  ^k{\hat{w_i}}$ is the weight of history number
$k$ before
going into collision number
$i$ and $ ^k{w_{i}^{*}}$ is the
weight of history number $k$ when coming out of collision number $i$.
In $R_t$ the sum is over line
integrals along the track of a test flight
from point of collision $ \vr_i$
to the next one at $\vr_{i+1} ,   i=0,1,
\dots , ^k{n-1} $, whereas in $R_c$ the sum is over
the points of collisions $\vr_i ,  i=1,2, \dots ,
^kn$. Here $i=0$ corresponds to the point of birth, and
$^kn$ the number of the last collision
along the track (i.e., the collision
which leads into the final state ``absorbed
particle").

The tallies are estimated as arithmetic means
\begin{equation}
R   =   \frac{1}{N}   \left ( \sum_{k=1}^{N}
(^k{R_t}  +  {^k{R_c}} ) \right )
\end{equation}
of the contributions of each test particle history k.

\subsection{Tracklength estimated volume tallies, UPTUSR}\label{sec3.2.1}

The tracklength estimators $R_t$ for detector function $g$ are
updated in subroutine UPDATE (defaults) and UPTUSR (...,WV,...)
(additional user supplied tallies), with
\begin{displaymath}
WV   =
{^k{w_{i}^{*}}}  /   | \vv | \quad .
\end{displaymath}
In EIRENE variables:
\begin{displaymath}
 WV = WEIGHT/VEL  \quad .
\end{displaymath}
Then, in order to estimate a response for a detector function $g$ on
additional tally number ITALV, a statement in subroutine UPTUSR
should read:
\begin{verbatim}
 DO IC=1,NCOU
   ICELL=NRCELL+NUPC(IC)*NR1P2+NBLCKA
   ADDV(ITALV,ICELL) = ADDV(ITALV,ICELL) + WV * CLPC(IC) *  g
 ENDDO
\end{verbatim}

Here $CLPC$ is the length $l$ of the flight in the cell ICELL, $g$
is assumed to be a constant in this cell in this example. In more
general cases, when detector $g$ varies along the track, the product
$l \cdot g$ is to be replaced by the line integral $\int g~ dl$
along this track.
\subsection{Collision estimated volume tallies, UPCUSR}\label{sec3.2.2}
The default collision estimated contributions $R_c$ are updated in
subroutine COLLIDE. Subroutine COLLIDE is called from the particle
tracing subroutines FOLNEUT and FOLION (defaults). Non default
collision tallies are updated in UPCUSR(...,WS,IND) (additional user
supplied tallies, called from COLLIDE). Here we have
\begin{displaymath}
WS   =  {^k{\hat{w_i}}}
  /   ( | \vv |  \Sigma ) \quad .
\end{displaymath}
In EIRENE variables:
\begin{displaymath}
WS = WEIGHT/SIGTOT = WEIGHT/VEL \cdot ZMFP \quad .
\end{displaymath}

$IND=1$ indicates a call before sampling from the collision
kernel, and $IND=2$ indicates calls with the parameters (velocity,
weight, ...) of the particle emerging from the collision. At each
collision, the subroutine UPCUSR is called once before and once
after the collision.

Then, in order to estimate a response for a detector function g on
additional tally number ICOLV, a statement in subroutine UPCUSR
should read:

\begin{verbatim}
 COLV(ICOLV,ICELL) = COLV(ICOLV,ICELL) + WS * g  .
\end{verbatim}

Here $g$ is evaluated at the point of a collision, which is known
to have taken place in cell ICELL in this call to subroutine
UPCUSR.

We note that track-length estimates ADDV reduce to collision
estimates COLV, if ADDV is updated at the points of collision only
and if the EIRENE variable ZMFP (local mean free path length) is
used rather than CLPD. This fact is used by EIRENE for ``condensed
particle species" (NFOL\$ flag in block 4). This is used, for
example, if the motion of test ions along field lines is not
followed explicitly but if, instead, the particles undergo a next
collision immediately at their places of birth (i.e.: quasi steady
state approximation ``QSSA" for this species).

\subsection{Snapshot estimated volume averaged tallies, UPNUSR}\label{sec3.2.3}

By default snapshot tallies are only scored in time dependent mode, see input flags in Section \ref{sec2.13}.
They can, however, also be used as a third alternative to tracklength and collision
estimators for any stationary response. In particular for diffusive processes, as e.g.
ion transport with Fokker Planck type Coulomb collision operators,
such ``snapshot averaging" is the most frequent option used in Monte Carlo applications to estimate
stationary moments. See paragraph B below.
\subsubsection{A: time dependent estimates}
The default
snapshot estimated tallies at a fixed time $t^*$, see Section
\ref{sec2.13}, are updated in subroutine TIMCOL. Currently these are
only the census tallies and the particle- and energy fluxes at
census. Subroutine TIMCOL is called from the particle tracing
subroutines FOLNEUT and FOLION (defaults), in case of time dependent
mode (see input block 13). Non default snapshot tallies can be
scored in UPNUSR (additional user supplied tallies, called from
TIMCOL). Here we have
\begin{displaymath}
WSNAP   =  {^k{\hat{w_i}}}   \quad .
\end{displaymath}
In EIRENE variables:
\begin{displaymath}
WSNAP = WEIGHT  \quad .
\end{displaymath}
Then, in order to estimate a response for a detector function $g$ on
snapshot tally number ISNV, a statement in subroutine UPNUSR should
read:
\begin{verbatim}
 SNAPV(ISNV,ICELL) = SNAPV(ISNV,ICELL) + WSNAP * g(...)  .
\end{verbatim}
E.g., with $g=1$ this is an unbiased estimate of cell averaged particle density at
time $t=t^*$, see (\ref{est}), i.e., just counting the particle weights at $t=t^*$
provides an unbiased estimate of the particle density.
This is intuitively clear, but also mathematically correct since
these snapshot estimators are formally derived as special cases of tracklength- collision
or surface flux estimators (interpreting the time-horizon sub-manifold as ``time-surface", by abuse of language)
by including in $g$ a delta function in time : $g(t,\vr,\vv)=\delta(t^*)\tilde g(\vr,\vv)$.

Subroutine UPNUSR must at least include module COMPRT (for the
particle weight WEIGHT). Here $g$ is evaluated at the point of time $t^*$
at which TIMCOL is called. Scaling of snapshot tallies is done by
interpreting the linear source strength ``FLUX" scaling factor
(given in input block \ref{sec2.7} for each stratum) as a total
number of particles, rather than as flux (particles/s), i.e. ``FLUX is the
integral of the source distribution not only over physical space, but also over time interval $DTIMV$ from initial time $t_0$
to time $t^*$ (Section
\ref{sec2.13}).
\subsubsection{B: stationary snapshot tallies}\label{sec3.2.3b}
A stationary value of snapshot tallies, to be compared with other
tallies in stationary mode, is obtained by setting the flag NTMSTP
of block \ref{sec2.13} to negative values. Stationary results in
this ``time dependent mode" are obtained by launching all test
flights at time $t_0$, but scoring snapshot tallies not only at
$t^*=t_1=t_0+DTIMV$, but also at $t_2=t_0+2~ DTIMV, t_3=t_0+3~
DTIMV, ...,$ i.e. at all times $t_n=t_0+n~ DTIMV$ until the flight
terminates. If the source distribution and the background medium
parameters are constant in time, then the snapshot score
contribution from time $t_i$ can be interpreted as score at
$t^*=t_1$ resulting from the source at earlier time $t_0-(i-1)
DTIMV$. Hence this stationary snapshot score at $t^*$ is in fact an
integration over all contributions from earlier times $t < t_0$.

When scaling the time-
snapshot tallies at $t=t^*$ in this stationary mode to absolute units, then an extra multiplicative factor $DTIMV$
(i.e., the time interval between subsequent scores along a particle history) is
applied for snapshot tallies, because the stationary (time-averaged) volume averaged
tallies are scaled interpreting the stratum source strength $FLUX$
as flux: ``atomic particles per second", whereas the snapshot
tallies require this same scaling factor $FLUX$ to be the absolute
number of (atomic) particles (flux integrated over one time interval
from $t_0$ to $t_1=t_0+DTIMV$).

Further considerations regarding the scaling factors for volume averaged
tallies in EIRENE, with respect to dimensionality of phase space of the problem
(1D, 2D, 3D, stationary or time-dependent) are given in Section \ref{sec1.2.0}.

\subsection{Surface averaged tallies, UPSUSR}\label{sec3.2.4}
\medskip
Finally we note that ``surface averaged tallies" (table \ref{tab3})
may be considered just as special cases of the ``volume averaged
tallies" described above, if appropriate use is made of $\delta$
-functions. Let therefore $\rho$ be a co-ordinate normal to a
surface S at the strike point $\underline{r}_S $ of a test flight
from $\underline{r}_i$ to $\underline{r}_{i+1}$ with speed unit
vector $\underline{\Omega} = \underline{v}  / (| \underline{v} |)$.
I.e., the surface is described locally by the equation $\rho=0$ at
the point $\vr_S \in S$, at which we may define an orthonormal basis
$ \ve_{\rho}, \ve_ {\rho'} , \ve_{\rho''}$. Furthermore, $\vr_S =
\vr_{i+1} $, which means that one may treat surface events exactly
in the same way as collisions with the background medium in our
terminologies. Then a ``surface response function" $ g_S $ defined
as

\[
g_S (j, \vr , \vv )   =
\delta(\rho)  \cdot  g(j, \vr , \vv )
\]

is the detector function for the same response $ R_g $ as previously
detector function $g$ was, but now surface averaged instead of
volume averaged:

\[
\int_{surface}   d^2 s  \ g  \cdot    f    =
\int_{volume}  d^3 r \
g_S   \cdot   f
\]

Therefore, the line integrals $ I_{l,i} $ in the formula given above
for track-length estimates, equation (\ref{tracklest}), now reduce
to:

\begin{equation}
I_{l,i}   =
\int_{ \vr_i}^{\vr_{i+1}}
dl \  \delta(\rho)  \cdot  g^{*} (l)   =
g^{*} (\rho=0)  /  \cos( \ve_{\rho}, \vOmega )
\end{equation}

if the surface is intersected
during the test-flight from $ \vr_i $
to $ \vr_{i+1} $ , and $ I_{l,i} = 0 $ else.
This can easily be seen by writing
\begin{eqnarray}
 \vr_{l,i} & = & \vr_i + l  \cdot
 \left(\begin{array} {c} \cos(\vOmega ,
\ve_{\rho}
)
\\  \cos(\vOmega , \ve_{ \rho '} )
\\  \cos(\vOmega,\ve_{ \rho ''}) \end{array}\right)
\end{eqnarray}
i.e., $ \rho= r_{\rho , i} + l  \cdot  \cos( \vOmega , \ve_{\rho} )
$ . In other words: tracklength estimators become formally identical
to collision estimators, for detector functions containing surface
collision delta functions.




Consequently, a statement in subroutine UPSUSR(WEIGHT,IND)
for this surface averaged tally must read:
\begin{displaymath}
ADDS(ITALS,IS)=ADDS(ITALS,IS) + WC   \cdot   g  \quad ;
\end{displaymath}

here $IS$ is the index of the surface which is being crossed,
$ITALS$ is the labelling number of the surface tally. Furthermore,
\begin{displaymath}
WC= {^k{\omega^{*}_i}}  /  [ |\vv |
  \cdot   \cos( \vOmega , \ve_{\rho} ) ] ,
\end{displaymath}
the detector function $g$ is evaluated at the strike point $ \vr_S $
and, as above, $ {^k{\omega^{*}_i}}  $ is the weight of the history
k after event i and before event i+1. Note:
\begin{displaymath}
 WEIGHT = {^k {w^{*}_i}}  =
^k {\hat{w_{i+1}}}
\end{displaymath}
in our terminology.

The flag IND has the value 1, if the particle is incident onto the
surface, and IND = 2 if the particle is emitted from the surface.

Some care is needed because the local surface normal unit vector $
\ve_{\rho} $ at the strike point $ \vr_S $ is not known in
subroutine UPSUSR unless the surface input variable ILIIN (see input
block 3) is positive (equal to 1,2,3 or 4) for the surface $S$. In
all other cases, this vector (defined by its Cartesian components
CRTX,CRTY,CRTZ in EIRENE) has to be computed in subroutine UPSUSR at
each entry, if it is needed for updating the surface averaged tally.

Subroutine UPSUSR is called whenever the default surface tallies are
updated. For one sided surface tallies (ILIIN=-2 or ILIIN=-4), therefore,
UPSUSR is also only called for one sided tallies only.

In addition to these calls, UPSUSR is also called from surface source routines.
For test particles emitted from source surfaces (or generated from incident bulk ions after
reflection or sputtering)
the call is UPSUSR(WEIGHT,2). For bulk ions (ITYP=4) incident onto a source surface
the call is UPSUSR(-WEIGHT,1).

Therefore, the tallies ADDS include the direct source contribution, if the source is a surface
source, whereas the default surface tallies don't.
\newpage

\section{The user surface reflection model REFUSR} \label{sec3.3}

\medskip

{\bf General remarks}:

\medskip

This subroutine is called in the initialization phase of an EIRENE
run, at entry RF0USR, from subroutine REFLEC, and at entry SP0USR, from
subroutine SPUTER. At these entries the user supplied reflection models
and sputtering models (if any),
respectively, may be initialized.

Later, during the Monte
Carlo sampling phase, it is called via entry RF1USR, whenever a test
flight intersects a surface NLLI, for which the fast particle reflection
model flag ILREF(NLLI) has been set equal to ILREF(NLLI)=3.

Furthermore, if
for a certain surface NLLI a user specified ``sputter model" is activated
by the flag ILSPT(NLLI)=3,
then REFUSR is called at the entry SP1USR,
whenever this surface is intersected by a test flight.

\medskip

{\bf Format of subroutine}

\medskip

\begin{verbatim}
      SUBROUTINE REFUSR
C USER SUPPLIED SURFACE INTERACTION ROUTINE
C INPUT : CO-ORDINATES OF INCIDENT PARTICLE
C OUTPUT: CO-ORDINATES OF EMITTED PARTICLE
* CALL PARMMOD
     :
     :
C  INITIALIZE USER SUPPLIED REFLECTION MODEL
      ENTRY RF0USR
     :
C  DEFINE SPECIES DEPENDENT RECYCLING COEFFICIENT HERE
C     RECYCT(ISPZ,MSURF)=.....
     :
      RETURN
      ENTRY RF1USR (XMW,XCW,XMP,XCP,IGASF,IGAST,ZCOS,ZSIN,EXPI,
     .              RPROB,E0TERM,*,*,*,*)
C  RETURN 1: EIRENE STANDARD ANGULAR DISTRIBUTION (DEP. ON ``EXPI")
C  RETURN 2: THERMAL MOLECULE MODEL (DEP. ON ``IGAST", ``E0TERM")
C  RETURN 3: THERMAL ATOM MODEL (DEP. ON ``IGAST", ``E0TERM")
C  RETURN 4: ABSORB PARTICLE AT THIS SURFACE
     :
     :
      RETURN
C
C  INITIALIZE USER SUPPLIED SPUTTERING MODEL
      ENTRY SP0USR
     :
C  DEFINE SPECIES DEPENDENT PHYSICAL SPUTTERING COEFFICIENT HERE
C     RECYCS(ISPZ,MSURF)=.....
C  AND ALSO CHEMICAL SPUTTERING COEFFICIENT
C     RECYCC(ISPZ,MSURF)=.....
     :
      ENTRY SP1USR
     :
     :
      RETURN
      END
\end{verbatim}
\newpage

\section{The user source sampling routine SAMUSR} \label{sec3.4}

\medskip

{\bf General remarks}:

\medskip

This subroutine is called from the point source sampling routine
SAMPNT, the surface source sampling routine SAMSRF, or from the
volume sampling routine SAMVOL, if the flag SORLIM(N, ISTRA)
(input block 7) for the spatial distribution of birth points on
sub-stratum N for stratum ISTRA has a negative value.

This subroutine returns to EIRENE the (cartesian) coordiantes of
the birth point

X0,Y0,Z0

cell index information

IRUSR,IPUSR,ITUSR,IAUSR,IBUSR

as well as local plasma background data at this place of birth:

TIWL,TEWL,DIWL,VXWL,VYWL,VZWL,EFWL,SHWL,WEISPZ

The input parameters (EIRENE input block 7 for primary sources)

SORAD1,SORAD2,SORAD3,SORAD4,SORAD5,SORAD6

can be used in this problem specific routine.

Dimensions, Precision:



\medskip

{\bf Format of subroutine}

\medskip

\begin{verbatim}
      SUBROUTINE SAMUSR (ISR,X0,Y0,Z0,
     .                   SORAD1,SORAD2,SORAD3,SORAD4,SORAD5,SORAD6,
     .                   IRUSR,IPUSR,ITUSR,IAUSR,IBUSR,
     .                   TIWL,TEWL,DIWL,VXWL,VYWL,VZWL,EFWL,SHWL,
     .                   WEISPZ)
      USE PRECISION
      USE PARMMOD
           :
           :  ! MORE MODULES, IF NEEDED
           :
      IMPLICIT NONE
      INTEGER, INTENT(IN) :: ISR, ISTR
      REAL(DP), INTENT(IN)   :: SORAD1,SORAD2,SORAD3,SORAD4,SORAD5,
     .                          SORAD6
      REAL(DP), INTENT(OUT) ::
     .          X0,Y0,Z0
      INTEGER,  INTENT(OUT) :: IRUSR,IPUSR,ITUSR,IAUSR,IBUSR
!   NPLS: NUMBER OF BULK ION SPECIES IN EIRENE INPUT BLOCK 5
!   NSPZ: TOTAL NUMBER OF PARTICLE SPECIES:
!         NSPZ=NPHOT+NATM+NMOL+NION+NPLS,
!         SEE INPUT BLOCKS 4 AND 5.
      REAL(DP), INTENT(OUT) ::
     .          TEWL,SHWL,
     .          TIWL(NPLS),DIWL(NPLS),
     .          VXWL(NPLS),VYWL(NPLS),VZWL(NPLS),
     .          EFWL(NPLS),WEISPZ(NSPZ)


!  INITIALIZE SAMPLING FOR SUBSTRATUM  ISR OF STRATUM ISTR
!  ISR:  NUMBER OF SUBSTRATUM (SEE INPUT BLOCK 7)
!  ISTR: NUMBER OF STRATUM
      ENTRY SM0USR (ISR,istr,sorad1,sorad2,sorad3,
     .                       sorad4,sorad5,sorad6)

      RETURN

      ENTRY SM1USR (ISR,X0,Y0,Z0,
     .              SORAD1,SORAD2,SORAD3,SORAD4,SORAD5,SORAD6,
     .              IRUSR,IPUSR,ITUSR,IAUSR,IBUSR,
     .              TIWL,TEWL,DIWL,VXWL,VYWL,VZWL,EFWL,SHWL,WEISPZ)
!  FIND BIRTH POINT COORDINATES. SUBSTRATUM ISR IF STRATUM ISTRA
!  HAS ALREADY BEEN IDENTIFIED. ISTRA (IF NEEDED) IS IN MODULE COMPRT.
!  INITIALIZE THOSE VARIABLES WHICH ARE NOT SET BELOW:
      X0  =0._DP
      Y0  =0._DP
      Z0  =0._DP

      IRUSR =0
      IPUSR =0
      ITUSR =1
      IAUSR =0
      IBUSR =1

      TEWL=0._DP
      SHWL=0._DP
      TIWL=0._DP
      DIWL=0._DP
      VXWL=0._DP
      VYWL=0._DP
      VZWL=0._DP
      EFWL=0._DP
      WEISPZ=0._DP

!  HERE COMES THE PROBLEM SPECIFIC DEFINITION OF BIRTH POINT SAMPLING


      RETURN

      END
\end{verbatim}
Note: The parameters EFWL, SHWL have been introduced in 2005.

Note: Correction in May 2006: In previous versions of this manual the parameters TEWL and TIWL have
been interchanged with respect to the calling statements in the EIRENE code. \\
(Thanks to Jose Guasp, (CIEMAT, Spain) for pointing this out.)

\newpage

\section{The user routines to overrule input data}
\label{sec3.5}
\subsection{The user geometry data routine GEOUSR} \label{sec3.5.1}
 to be written
\newpage
\subsection{User supplied background data routine PLAUSR}
\label{sec3.5.2}
 to be written
\section{The user routines for profiles PROUSR}
\label{sec3.6}
\medskip

{\bf General remarks}:

\medskip

Purpose: set input tallies (background medium), on the grid (one number per tally per cell)
For more general input tallies, evaluated ``on the fly" at a particular position, see routine VECUSR
below (\ref{sec3.6.2}).

This subroutine is called from the subroutine PLASMA
if the input flag INDPRO (input block 5) has the value INDPRO(IPRO)=5
for a particular background tally IPRO.
Only the data for cells from the standard mesh (ICELL=1,..,NSURF)
are used. Default ``vacuum data" are set for all cells in the ``additional
cell region" ICELL=NSURF+1,...,NSBOX.

PROUSR is called, furthermore, from subroutine INPUT to provide cell volumes for all
cells (ICELL=1,...,NSBOX), if the input flag INDPRO(12) has the value INDPRO(12)=5

\medskip

{\bf Format of subroutine}

\medskip

\begin{verbatim}
      SUBROUTINE PROUSR (RHO,INDEX,P0,P1,P2,P3,P4,P5,PVAC,NDAT)
* CALL PARMMOD
     :
     :
      IF (INDEX.EQ.0) THEN
C  DATA FOR TE-PROFILE, NDAT=NSURF
        DO ICELL=1,NDAT
          RHO(ICELL)=......
        ENDDO
      ELSEIF (INDEX.EQ.....) THEN
     :
     :
      ELSEIF (INDEX.EQ.5+5*NPLS) THEN
C  DATA FOR VOL, NDAT=NSBOX
     :
     :
      ELSE
        WRITE (6,*) 'INVALID INDEX IN PROUSR '
        CALL EXIT
      ENDIF
      RETURN
      END
\end{verbatim}

Depending upon the value of INDEX, one has to specify a background tally
on the array RHO, RHO(J),J=1,NDAT.
These are:
\begin{itemize}
\item
INDEX = 0     : \\
Electron temperature (eV)
\item
INDEX = 1     :\\
Ion temperature (eV), NPLSI calls, one for each
ion species. I.e.:
\begin{verbatim}
DO IPLS=1,NPLSI
  CALL PROUSR(HELP,1+0*NPLSI,TI0(IPLS),...,NSURF)
  CALL RESETP(TIIN,HELP,IPLS,1,NPLS,NSURF)
ENDDO
\end{verbatim}
\item
INDEX = 1+ NPLS     : \\
Ion density ($cm^{-3}$), NPLSI calls, one for each
ion species
\item
INDEX = 1+2*NPLS    :   \\
Ion drift velocity, x direction (cm/s), NPLSI calls, one for each
ion species
\item
INDEX = 1+3*NPLS    :   \\
Ion drift velocity, y direction (cm/s), NPLSI calls, one for each
ion species
\item
INDEX = 1+4*NPLS    :   \\
Ion drift velocity, z direction (cm/s), NPLSI calls, one for each
ion species
\item
INDEX = 1+5*NPLS    :   \\
Magnetic field, x component, (AU)
\item
INDEX = 2+5*NPLS    :   \\
Magnetic field, y component, (AU)
\item
INDEX = 3+5*NPLS    :   \\
Magnetic field, z component, (AU)
\item
INDEX = 4+5*NPLS    :   \\
Magnetic field, magnitude, (T)
\item
INDEX = 5+5*NPLS    :   \\
Cell volume, ($cm^3$)
\item
INDEX = 6+5*NPLS    :   \\
Additional tally, NAINI calls, one for each additional quantity.
\end{itemize}

\subsection{User supplied background data routine VECUSR}
\label{sec3.6.2}
This routine is called at a particular position (X0,Y0,Z0) in grid cell NCELL, and returns a vector quantity.
Position and cell number information are taken from module EIRMOD\_COMPRT.F.
It is used in particular if the magnetic field, or plasma flow field, are needed with a higher spatial
resolution than the one provided by the standard input plasma (background) tallies, which are cell averaged,
i.e. constant within a cell.

\begin{verbatim}
      SUBROUTINE VECUSR(INDEX,VEC_X,VEC_Y,VEC_Z,IPLS)
\end{verbatim}
This routine is called in case INDPRO(5)=8 for the B-field tally (e.g. from routine BFIELD.F), and in case (INDPRO(4)=8)
for the background flow velocity tallies (from many places in the code),
to provide either a local magnetic field (unit) vector, or a local plasma (background medium)
flow velocity vector (cm/s), respectively, in cartesian coordinates.

Current implementation:
\begin{itemize}
\item
INDEX = 1:   \\
return the magnetic field unit vector, cartesian coordinates. This local (at position X0,Y0,Z0, in cell NCELL) vector is used
instead of cell averaged input tally no. 5:

[BXIN(ncell),BYIN(ncell),BZIN(ncell)].

VECUSR with INDEX=1 is called from routine BFIELD.F
\item
INDEX = 2:   \\
return plasma (background) flow velocity vector, cartesian coordinates, for EIRENE background species IPLS.
This local (at position X0,Y0,Z0, in cell NCELL) vector is used
instead of cell averaged input tally no. 4:

[VXIN(IPLS,ncell),VYIN(IPLS,ncell),VZIN(IPLS,ncell)].

Note: species index IPLS is the EIRENE background species index, not to be confused with the mapped species index MPLSV(IPLS).
\end{itemize}

\section{User supplied post-processed tally routine TALUSR}
\label{sec3.7}
to be written
\begin{verbatim}
      SUBROUTINE TALUSR(ICOUNT,VECTOR,TALTOT,TALAV,
     .                  TXTTL,TXTSP,TXTUN,ILAST,*)
\end{verbatim}
\section{User supplied ``general geometry block"}
\label{sec3.8} The following routines have to be provided for the
general geometry options (LEVGEO=10 option, NLGEN=TRUE), in which no
specific geometry data are available from EIRENE input, and all the
geometrical parameters for particle tracing and scoring of tallies
are transferred from outside.

\begin{itemize}
\item
INIUSR (initialize user specified geometry block), see \ref{sec3.8.0}
\item
LEAUSR(x) (return cell number, for any given position x), see
\ref{sec3.8.1}
\item
TIMUSR (flight time to next cell boundary, and next cell number
or surface number), see \ref{sec3.8.2}
\item
VOLUSR (volume of each grid cell),
see \ref{sec3.8.3}
\item
NORUSR (outer surface normal, for any given position on a surface),
see \ref{sec3.8.4}
\end{itemize}
In case NLGEN, only a reduced set of the standard EIRENE options is available, and the
five user routines mentioned above and described below may have to be
supplemented by some further options from
the problem specific segment USER.F .
\begin{list}{ }{}
\item[{\bf Input block 5 }]

only the INDPRO = 3, 5 and 6 options are available. INDPRO = 3 provides
constant profiles, same value in each cell.  Hence, in case of
non-constant
background parameters one has to resort to subroutine PROUSR
(INDPRO=5), or to the code segment INFCOP (INDPRO=6) for coupling to another
source (code) for the data of the background medium.
\item[{\bf Input block 7 }]
only the point source option is available. In case of surface
sources or volume sources, the subroutine SAMUSR has to be called,
see \ref{sec3.4}.
\item[{\bf Input block 11 }]
only the printout options are available, no graphics output
options are available. However, subroutine PLTUSR(LOG,J) is called
from the 2D geometry plotting routine PLT2D (LOG=.TRUE., J = number of
surface to be plotted) and
from the 3D geometry plotting routine PLT3D (LOG=.FALSE., J =
number of surface to be plotted)
\end{list}
\subsection{Subroutine INIUSR}
\label{sec3.8.0} This subroutine is called from subroutine INPUT,
after reading from the formatted input file and after (optional)
calls to interfacing routines INFCOP. Any initialization
(including definition of COMMON blocks) for the user geometry
package may be done here.
\subsection{Subroutine LEAUSR}
\label{sec3.8.1}

Identify cell index from a point given by its three cartesian coordinates.

The function LEAUSR is called from functions LEARC1
and LEARCA in case of LEVGEO=10. The call is
\begin{verbatim}
NCELL=LEAUSR(X,Y,Z)
\end{verbatim}
Here X,Y,Z are the Cartesian
coordinates of a point inside the computational domain in centimeters.
LEAUSR must then return
the number of the cell to which this point belongs.

Note (see section \ref{sec2.2}): Any particular cell in EIRENE can be specified in
one of two possible ways:

Either by giving the 5 cell numbers:

NRCELL, NPCELL, NTCELL, NBLOCK, NACELL

in the first,(radial) second (poloidal) third (toroidal) grid or additional cell region,
or, alternatively, by giving the cell index in the 1-dimensional cell array

NCELL

\medskip
The relation between these two options is:

NCELL = NRCELL + ((NPECLL-1)*NR1P2) + (NTCELL-1))*NP2T3 + NBLCKA

with

NBLCKA = (NBLOCK -1) * NSTRDT +NRADD

LEAUSR is expected to return the cell index NCELL.



\subsection{Subroutine TIMUSR}
\label{sec3.8.2}
The subroutine TIMUSR is called from TIMER (block GEO3D) in case of LEVGEO=10 (general external geometry option).
Distinct from the other geometry levels in TIMER, the number of the nearest intersected surface (MRSURF) along the trajectory
is not returned in general. Only if this nearest surface is one of the ``non-default standard surfaces" defined in input
block 3a, this information is returned (by a negative value of argument NEWCEL.

The call is
\begin{verbatim}
CALL TIMUSR (NRCELL,X0,Y0,Z0,VELX,VELY,VELZ,
             NJUMP,NEWCEL,TIM,ICOS,IERR,NPANU,NLSRFX)
\end{verbatim}
{\bf Input}

\begin{list}{ }{}
\item[{\bf NRCELL }]
Actual cell number, for which the intersection to the cell boundary has to be found.
If NJUMP=0, i.e., for the first call of a new trajectory (e.g., after a
collision), the starting point X0,Y0,Z0 must be in that cell.
In later calls for the same trajectory (NJUMP $>$ 0), NRCELL has been automatically updated,
i.e., it must not necessarily contain the starting point
X0,Y0,Z0, and the flight time (distance) from the starting point is accumulated.
\item[{\bf X0,Y0,Z0 }]
Cartesian coordinates of the starting point of a trajectory.
\item[{\bf VELX,VELY,VELZ }]~

unit (speed-) vector pointing in the
direction of the flight.
\item[{\bf NJUMP }]
\begin{list}{ }{}
\item[{\bf $= 0$ }]
this is the first call for a particular trajectory.
\item[{\bf $\neq 0$ }]
this is a later call for a particular trajectory. The initial position and
speed unit vector are the same as in the previous call. Hence: the geometrical
parameters depending only on those need not be evaluated again.
\end{list}
\item[{\bf NPANU }]
number of EIRENE test flight (particle history). Only for diagnostic printout from TIMUSR.
\item[{\bf NLSRFX }]
TRUE:  the starting position (X0,Y0,Z0) of the current trajectory is exactly on the cell boundary of cell
NRCELL.  This information may by used inside TIMUSR
to detect and exclude flight pathes of zero length (time)

FALSE:  Particle trajectory does not start at the cell boundary of cell NRCELL, but inside this cell.
\end{list}

{\bf Output}

\begin{list}{ }{}
\item[{\bf NEWCEL }]
\begin{list}{ }{}
\item[{\bf $< 0 $}]
trajectory has intersected one of the non-default surfaces (input block 3a).
ABS(NEWCEL) is the number of that surface (corresponding to running index
ISTSI in input block 3a).
\item[{\bf $> 0$ }]
number of next cell, in which the flight would continue, if no collision event
stops the track already earlier. I.e., cell number of the neighbor cell in the direction
of the flight.
\end{list}
\item[{\bf TIM }]
distance (cm) from [X0,Y0,Z0]
to the nearest intersection of the trajectory with a
boundary of cell NRCELL.  Since the speed vector [VELX,VELY,VELZ] is normalized to one (cm/s),
this distance is also the flight time (s)
\item[{\bf ICOS }]
only relevant, if the trajectory has intersected a non-default surface.
In this case, ICOS is the sign (+1 or -1) of the cosine of the angle of
incidence against the surface normal (this latter unit vector, [CRTX,CRTY,CRTZ]
may, e.g., be found in TIMUSR by a call
to NORUSR, see section \ref{sec3.8.4}).
\item[{\bf IERR }]
error flag. Presently
any value different from 0 will lead to an immediate end of the run.
See subroutine TIMER in code segment GEO3D.F


\end{list}
\subsection{Subroutine VOLUSR}
\label{sec3.8.3} The subroutine VOLUSR is called from subroutine
VOLUME in case of LEVGEO=10. The call is
\begin{verbatim}
CALL VOLUSR(N,A)
\end{verbatim}
Here N is the total number of cells in this run, and
A is an array of length N, containing the N cell volume values in $cm^3$.
\subsection{Subroutine NORUSR}
\label{sec3.8.4} The subroutine NORUSR provides the outer surface normal unit vector CX,CY,CZ at surfaces,
i.e. after a particle trajectory has intersected a ``non default grid surface" at point X,Y,Z.

It is called in the LEVGEO=10 option from subroutine
STDCOL,  after all surface intersection settings and possible switches have been carried out,
or directly via the entry point STDNOR of STDCOL. The call is
\begin{verbatim}
CALL NORUSR(M,X,Y,Z,CX,CY,CZ,SCOS,VELX,VELY,VELZ,NRCELL)
\end{verbatim}


Here [X,Y,Z] are the cartesian coordinates
of an intersetcion point, located on non-default standard surface no. M. The routine must return
the surface normal unit vector [CX,CY,CZ].

If surface M has a periodicity condition (i.e. surface flag ILIIN(M) $>$3 in input block 3a (\ref{sec2.3a})
then this routine must also carry out the periodicity transformation for the test particle (which may be different
for neutral particles and for charged (test ion) particles).

Periodicity is carried out, in general, by moving particles to another surface and/or altering its speed unit vector.

In this cases the code expects on output in [X,Y,Z] the new coordinates, on integer variable M the new surface number, and, if applicable,
on [VELX,VELY,VELZ] the new speed unit vector. NRCELL must contain the cell number, into which the new flight enters from surface M.

Currently NORUSR is called for all non default surfaces in the general (external) geometry level LEVGEO=10, i.e. for NRGEN = T.

It is also called (currently without warning or conditioning) for tetrahedral grids at periodicity surfaces.
